\uuid{UuN2}
\titre{Convergence}
\auteur{bodin ; hanani ; mzari}
\organisation{exo7}
\licence{CC BY-NC-SA 4.0}
\datecreate{2020-01-20}
\niveau{L1}
\module{Analyse}
\chapitre{Suite}
\sousChapitre{Autre}
\difficulte{2}
\theme{Suites, convergence}
\exo7qcmid{lille-1.suites.017}

\contenu{
\begin{question}
\qclassification{121.00}
\qtags{motcle=convergence}
Soit $\displaystyle u_n=\frac{1}{1^2}+\frac{1}{2^2}+\frac{1}{3^2}+\dots+\frac{1}{n^2}$ et $\displaystyle v_n=u_n+\frac{1}{n}$. Quelles sont les bonnes réponses ?
\begin{answers}
    \bad{Si les limites existent, alors $\displaystyle \lim _{n\to +\infty}u_n<\lim _{n\to +\infty}v_n$.}
    \bad{Les suites $(u_n)$ et $(v_n)$ sont divergentes.}
    \good{Les suites $(u_n)$ et $(v_n)$ sont adjacentes.}
    \good{Les suites $(u_n)$ et $(v_n)$ convergent vers la même limite finie.}
\end{answers}
\begin{explanations}
On vérifie que $(u_n)$ est croissante, $(v_n)$ est décroissante et que
$$\displaystyle \lim _{n\to +\infty}(u_n-v_n)=0.$$
Donc $(u_n)$ et $(v_n)$ sont adjacentes. En conséquence, elles convergent vers la même limite finie.
\end{explanations}
\end{question}
}
